% ============================================================================
%  COMS2 2026 submission — Springer CCIS (Communications in Computer and
%  Information Science). CCIS uses the Springer LNCS class.
%
%  BUILD
%    1. Download the Springer LNCS LaTeX package (llncs.cls, splncs04.bst):
%       https://www.springernature.com/gp/computer-science/lncs/conference-proceedings-guidelines
%    2. Place llncs.cls beside this file.
%    3. pdflatex coms2_paper   (run twice so cross-references resolve)
%
%  BEFORE SUBMITTING — replace every <...> placeholder:
%    - author names, affiliations, e-mails
%    - ORCID identifiers (mandatory per the call for papers)
%  Then verify every reference against its DOI. Do not submit a citation you
%  have not personally opened.
% ============================================================================

\documentclass[runningheads]{llncs}

\usepackage[T1]{fontenc}
\usepackage[utf8]{inputenc}
\usepackage{graphicx}
\usepackage{booktabs}
\usepackage{amsmath}
\usepackage{xcolor}
\usepackage{tikz}
\usetikzlibrary{arrows.meta, positioning}
\usepackage[hidelinks]{hyperref}

\begin{document}

% ----------------------------------------------------------------------------
\title{Detecting and Migrating Quantum-Vulnerable Cryptography in Blockchain
Source Code: A Precision-Oriented Toolkit with Blockchain-Anchored Audit}

\titlerunning{Detecting Quantum-Vulnerable Cryptography in Blockchain Code}

% ---------------------------------------------------------------------------
% AUTHOR BLOCK
%
% COMPLETE. Name, ORCID, affiliation and e-mail are all filled in.
% ORCID 0009-0003-1356-7903 -- ISO 7064 check digit verified.
%
% Optional improvements, neither required:
%   - an institutional e-mail in place of the personal one, which
%     corroborates the affiliation rather than asking a reviewer to take it
%     on trust
%   - supervisor co-authorship, normally expected for MSc project work;
%     see the template below
%
%   AFFILIATION    Amity Institute of Information Technology (AIIT), Amity
%                  University Rajasthan, Jaipur -- as stated by the author.
%                  The degree programme (MSc Cyber Security) does not belong
%                  in an LNCS affiliation line; it goes in the biography or
%                  acknowledgements if anywhere.
%
%   SUPERVISOR     This work is final-year MSc project work. Your guide
%                  almost certainly expects to be a co-author, and a
%                  supervisor on the author line also gives a reviewer
%                  someone accountable for the claims. Settle this BEFORE
%                  submitting -- adding an author after acceptance is
%                  awkward, and omitting a guide who expected credit is
%                  worse. See the co-author template below.
%
%   E-MAIL         Filled with benidictdavid18@gmail.com, which matches the
%                  author name spelling. Replace it with your Amity address
%                  (@s.amity.edu or similar) if you have one: an institutional
%                  address corroborates the affiliation, whereas a personal
%                  one leaves a reviewer taking it on trust. Accepted either
%                  way, but the institutional address is the stronger choice.
%
% ADDING CO-AUTHORS: supervisors and contributors go here too. For a second
% author at the SAME institution:
%
%   \author{Benidict David\inst{1}\orcidID{...} \and
%           Co Author\inst{1}\orcidID{...}}
%   \authorrunning{B. David and C. Author}
%
% For a co-author at a DIFFERENT institution, give them \inst{2} and add a
% second block to \institute separated by \and.
% ---------------------------------------------------------------------------

\author{Benidict David\inst{1}\orcidID{0009-0003-1356-7903}}

\authorrunning{B. David}

\institute{%
  Amity Institute of Information Technology,\\
  Amity University Rajasthan, Jaipur, Rajasthan, India\\
  \email{benidictdavid18@gmail.com}%
}

\maketitle

% ----------------------------------------------------------------------------
\begin{abstract}
Blockchains rely on RSA and ECDSA, and a large enough quantum computer breaks
both. They suffer more than most systems from this, because every public key
they have recorded stays public forever, so an attacker who gets a quantum
computer years from now can still attack today's data. Before a company can move
to safer algorithms it needs a list of exactly where in its code the weak ones
are used. Existing smart contract tools do not produce that list: they look for
logic bugs, not for which cryptography a program uses. We present a tool that
finds these places in Solidity, Python, JavaScript and Go code, gives the
NIST-approved replacement for each, writes a report for regulators, and records
a tamper-proof summary on a blockchain. Our main result concerns
accuracy. A text-matching tool of this kind is easy to make find everything but
hard to stop from being wrong, and two causes explain almost all its mistakes:
it matches algorithm names written in comments, and it reports one line several
times when several patterns match it. We fix both, and replace a risk score that
gave nearly every file the same value. On a test set of seven files with
nineteen marked locations, three written to trick the tool, these fixes raise
precision from 0.633 to 1.000 while still finding every real case. Because we
wrote that test set ourselves, we explain what it does and does not prove. The
tool also runs in a browser, giving identical results, so private code never
leaves the machine.

\keywords{Post-quantum cryptography \and Blockchain security \and Static
analysis \and ML-DSA \and ML-KEM \and Cryptographic bill of materials \and
Regulatory compliance.}
\end{abstract}

% ============================================================================
\section{Introduction}
\label{sec:intro}

Blockchains depend on only a few cryptographic algorithms. Ethereum and Bitcoin
sign transactions with ECDSA on the secp256k1 curve. Smart contracts check
signatures using the \texttt{ecrecover} function. The servers and scripts around
a blockchain usually use RSA to exchange keys and SHA-256 to make hashes. All of
these depend on maths problems that a quantum computer can solve quickly.
Shor's algorithm~\cite{shor1997} can factor large numbers and solve discrete
logarithms fast, which breaks RSA and ECDSA completely. Grover's
algorithm~\cite{grover1996} is a weaker attack: it makes searching twice as
fast, which cuts the strength of hashes and block ciphers in half. AES-128 drops
to about 64 bits of security and SHA-256 to about 128 bits.

Blockchains have a problem here that most systems do not. Normally, when you
stop using a weak algorithm, the old data is gone and an attacker cannot reach
it. A public blockchain works the other way round. Every public key and every
signature it has ever stored is still there and still public. So an attacker who
gets a quantum computer in ten years can still attack everything written today.
Mosca~\cite{mosca2018} explains the timing simply: if the time your data must
stay secret, plus the time your migration takes, is longer than the time until
quantum computers arrive, you have already left it too late. For a blockchain
the data must stay secret forever, so by this rule the work should already have
started.

The algorithms to move to are now settled, so there is no technical excuse to
wait. In August 2024 NIST published three final standards: FIPS 203 for ML-KEM
(key exchange)~\cite{fips203}, FIPS 204 for ML-DSA
(signatures)~\cite{fips204} and FIPS 205 for SLH-DSA (hash-based
signatures)~\cite{fips205}. NIST has also published a timetable~\cite{nistir8547}
for when the old algorithms stop being allowed. In India, the Reserve Bank of
India's cybersecurity framework and CERT-In's directions already require
regulated companies to know what cryptography they are using.

The step nobody has solved is the first one. Before a company can change
anything, it has to know where the weak algorithms are used in its code, and
know it precisely enough to hand the work to a developer. This is a listing
problem, and it is not what existing smart contract tools do.
Slither~\cite{slither}, Mythril and Manticore look for bugs such as reentrancy,
number overflow and missing access checks. They do not report which
cryptographic algorithms the code uses, because using ECDSA has never been a bug
until now. CryptoGuard~\cite{cryptoguard} does look at cryptography, but it
checks for wrong usage in Java programs, not for quantum risk in blockchain
code.

\subsection{Problem Statement}

The problem we solve is this. Given a source file from a blockchain project,
produce a list of the exact places where quantum-vulnerable cryptography is
used. The list has to be detailed enough to plan the migration, and accurate
enough that a developer reading it still trusts the tool.

Being accurate is the hard part, and it is easy to underestimate how hard. A
tool like this is normally built by searching the source code for patterns. It
is easy to make such a tool find every real case. It is much harder to stop it
reporting things that are not real. Two mistakes cause most of the trouble.

First, algorithm names appear a lot in comments and in text strings. This is
worst in exactly the code we care about, because a file that is being migrated
usually has comments explaining which old algorithms were removed. A tool that
searches the raw text reports those comments as problems.

Second, one line of code often matches several patterns at once, so the tool
reports the same place two or three times. The count then looks bigger than the
number of places a developer actually has to visit.

Both mistakes make every later number wrong too, including the risk score and
the estimate of how much work the migration needs.

\subsection{Contributions}

\begin{enumerate}
  \item \textbf{Two fixes that make the tool accurate.} We blank out comments
  and text strings before searching, but we replace each character with a space
  instead of deleting it, so the line numbers we report stay correct. We also
  group matches that land on the same line for the same algorithm, so one place
  in the code gives one result (Section~\ref{sec:method}).

  \item \textbf{A risk score that actually varies.} The earlier version of our
  tool added up severity points and capped the total at 100. Four serious
  findings already reach the cap, so almost every file scored 100. Across the
  scans we collected while building the tool, the average was 99.5. Our new
  score combines three things instead: how serious the worst finding is, how
  many different algorithms are involved, and how many findings there are per
  hundred lines (Section~\ref{sec:scoring}).

  \item \textbf{A test set with deliberately tricky files.} Seven files with
  nineteen marked locations. Three of the files contain no real problems at all:
  one already uses post-quantum algorithms, one only mentions the old algorithms
  in comments, and one has names like \texttt{bandwidth} and
  \texttt{sha256\_cache} that look cryptographic but are not. We release the
  files and the marking rules (Section~\ref{sec:benchmark}).

  \item \textbf{A test showing which fix does what.} We run the tool in three
  versions to show how much each fix helps, and we explain clearly why our
  scores look better than the tool would do on real code
  (Section~\ref{sec:results}).

  \item \textbf{The same tool in two languages, checked to agree.} We wrote the
  detection logic in Python and again in JavaScript, and we check that both give
  exactly the same results on every test file. The JavaScript version runs in a
  web browser, so a company can check its private code without uploading it
  anywhere (Section~\ref{sec:dualruntime}).

  \item \textbf{Output that companies can actually use.} A machine-readable
  list of cryptography in the standard CycloneDX format, reports written for
  regulators, and a record saved on a blockchain that proves when the check was
  done (Sections~\ref{sec:cbom} and~\ref{sec:anchor}).
\end{enumerate}

% ============================================================================
\section{Background and Related Work}
\label{sec:background}

\subsection{What Quantum Computers Break}

Today's public-key cryptography works because some maths problems are too slow
for normal computers to solve. RSA is safe because factoring a large number
into its primes takes too long. ECDSA is safe because working backwards from a
point on an elliptic curve takes too long. Diffie--Hellman relies on a similar
problem. Shor's algorithm~\cite{shor1997} solves all three quickly on a quantum
computer. These algorithms are not weakened by it; they stop working
altogether.

The only open question is when the hardware arrives. Webber et
al.~\cite{webber2022} estimated that stealing a Bitcoin private key within an
hour would need hundreds of millions of physical qubits, far beyond anything
built so far, but such estimates keep falling as error correction improves, and
in practice the deadline companies work to is set by regulators rather than by
any single estimate. Grover's algorithm~\cite{grover1996} is a milder threat: it
speeds up brute-force search, halving the strength of hashes and block ciphers.
The fix there is bigger sizes, AES-256 and SHA-384 or SHA-512, rather than a
different kind of algorithm.

\subsection{What to Replace Them With}

Table~\ref{tab:nist} lists the three standards NIST finalised in August 2024.
Bernstein and Lange~\cite{bernstein2017} give a good overview of the maths
behind them.

\begin{table}[t]
\centering
\caption{The finalised NIST post-quantum standards and what each one replaces.}
\label{tab:nist}
\begin{tabular}{@{}llll@{}}
\toprule
\textbf{Standard} & \textbf{Algorithm} & \textbf{Based on} & \textbf{Replaces} \\
\midrule
FIPS 203 & ML-KEM (Kyber)      & Lattices       & RSA key transport, DH, ECDH \\
FIPS 204 & ML-DSA (Dilithium)  & Lattices       & RSA signatures, ECDSA \\
FIPS 205 & SLH-DSA (SPHINCS+)  & Hash functions & ECDSA, long-lived keys \\
\bottomrule
\end{tabular}
\end{table}

One practical problem shapes what we can recommend for code that runs on-chain.
ML-DSA signatures and keys are much bigger than ECDSA ones, and running lattice
verification inside a smart contract costs more gas than is realistic today. So
a contract that checks signatures on-chain cannot simply swap one function for
another. It has to be restructured: the signature is checked off-chain, and the
contract only records that the check passed. Our tool says this plainly instead
of suggesting a replacement that will not run.

There is a second limit. Ethereum derives wallet addresses using Keccak-256, so
a deployed contract cannot switch to a different hash without breaking every
address that depends on it. We therefore only recommend changing hashes in
off-chain code.

\subsection{Why Existing Tools Do Not Help}

Tools for checking smart contracts are well developed. Slither~\cite{slither}
analyses Solidity code and finds a long list of known bug patterns. Mythril
runs the compiled contract symbolically, and Manticore can do both. None of
them reports which cryptographic algorithms a contract uses, because using
ECDSA has never counted as a bug in their view of the world.

Checking cryptography has been studied separately.
CryptoGuard~\cite{cryptoguard} finds mistakes such as hard-coded keys and
insecure modes in Java programs, and CrySL~\cite{crysl} lets you write rules
for how a crypto library should be used. Both look for \emph{wrong use of safe
algorithms}. Our problem is the opposite: \emph{correct use of algorithms that
are about to become unsafe}. A misuse checker reports nothing at all for a
textbook-perfect ECDSA implementation, which is exactly the code we need to
find.

The idea of listing the cryptography a system uses has recently been given a
standard form. CycloneDX~\cite{cyclonedx} added a
\texttt{cryptographic-asset} component type, which lets you describe an
algorithm, what kind of primitive it is, and where it runs, in a format other
tools can read. We produce this format rather than inventing our own, so our
output fits into software supply-chain tools companies already use.

As far as we know, no published tool combines all four of: finding quantum-weak
cryptography in blockchain code, planning the migration against the standards,
exporting a machine-readable inventory, and producing a verifiable record. So
part of our contribution is simply joining these up. The research contribution,
though, is the accuracy work in Section~\ref{sec:method}, which applies to any
tool that searches source code for cryptography.

% ============================================================================
\section{Threat Model and Design Goals}
\label{sec:threat}

It is worth stating exactly who we are defending against and what we do not try
to do, because that boundary explains several of our design choices and all of
our limitations.

\paragraph{The attacker.} We assume an attacker who follows a \emph{harvest
now, decrypt later} strategy. They record public blockchain data continuously
today, and at some unknown point in the future they get a quantum computer.
From then on they can work out private keys from the public keys they recorded
and forge signatures. Note that this attacker never needs to break into the
company's systems, because everything they need is already public.

\paragraph{What we do and do not check.} Our tool answers one question:
\emph{which lines of this code use an algorithm that the attacker above can
break?} It does not check whether the protocol design is sound, whether keys
are stored safely, whether random numbers are generated properly, or whether a
correctly-named algorithm is being used correctly. Those are separate jobs. Our
output is a to-do list for a migration, not a security audit, and we say so in
the generated reports as well as here.

\paragraph{What this means for the design.} Four goals follow.
\textbf{G1, be accurate rather than thorough}: a developer has to be able to
trust that every line we report really needs work, because a list that flags
comments as problems gets thrown away the first time someone checks it.
\textbf{G2, be specific}: every result needs an exact line number, a named
replacement, and the standard it comes from.
\textbf{G3, keep the code private}: checking private blockchain code must not
require uploading it, because the companies with the most to lose are the least
willing to send their source to someone else.
\textbf{G4, be verifiable}: when a company shows a report to a regulator or an
insurer, that party needs to be able to check it is genuine and when it was
made, without having to trust the company.

% ============================================================================
\section{System Architecture}
\label{sec:arch}

The tool works as five steps in a row (Fig.~\ref{fig:arch}). Each step is
separate and passes its results on as JSON, so any one of them can be replaced
without touching the others.

\begin{figure}[t]
\centering
\begin{tikzpicture}[
  node distance=3mm,
  box/.style={draw, rounded corners=2pt, align=center, font=\scriptsize,
              minimum height=8.5mm, inner xsep=2pt, text width=17mm},
  out/.style={draw, dashed, rounded corners=2pt, align=center,
              font=\scriptsize, minimum height=7mm, text width=19mm},
  ar/.style={-{Latex[length=1.5mm]}, thick}
]
\node[box] (src) {Source\\file};
\node[box, right=of src] (scan) {Hide comments,\\then search};
\node[box, right=of scan] (dedup) {Group by\\line};
\node[box, right=of dedup] (score) {Risk\\score};
\node[box, right=of score] (plan) {Migration\\plan};

\draw[ar] (src) -- (scan);
\draw[ar] (scan) -- (dedup);
\draw[ar] (dedup) -- (score);
\draw[ar] (score) -- (plan);

\node[out, below=7mm of dedup] (cbom) {CycloneDX\\list};
\node[out, below=7mm of score] (rep) {Report for\\regulator};
\node[out, below=7mm of plan] (anc) {Blockchain\\record};

\draw[ar, gray] (dedup) -- (cbom);
\draw[ar, gray] (score) -- (rep);
\draw[ar, gray] (plan) -- (anc);
\end{tikzpicture}
\caption{How the tool processes a file. Solid boxes are the analysis steps;
dashed boxes are the documents produced from the step above.}
\label{fig:arch}
\end{figure}

Steps 1 to 3 are the detection itself and are described in
Section~\ref{sec:method}: hide everything that is not real code, search what is
left, group repeated matches per line, and score the result. Names found only
inside comments are reported separately rather than discarded, so a reviewer can
see what the tool chose not to count.

Step 4 sorts the results by algorithm group and gives, for each group, the
replacement, the governing standard, the affected line numbers and a code
example in the right language. For Solidity it gives the off-chain pattern of
Section~\ref{sec:background}, because the direct replacement will not run. Step
5 produces the three documents: the CycloneDX list (Section~\ref{sec:cbom}), a
report in the relevant regulator's format, and the blockchain record
(Section~\ref{sec:anchor}).

\section{Detection Method}
\label{sec:method}

\subsection{What We Search For}
\label{sec:patterns}

We cover five groups of algorithms (Table~\ref{tab:families}). Each pattern
looks for a function call or an import rather than just the algorithm's name.
That helps a little, but it is the weakest of our accuracy measures; the two
described after it do the real work.

\begin{table}[t]
\centering
\caption{The five algorithm groups, how serious each one is, and what it should
become. HIGH means the attacker breaks it completely; MEDIUM means it is only
weakened.}
\label{tab:families}
\begin{tabular}{@{}lclp{32mm}@{}}
\toprule
\textbf{Group} & \textbf{Level} & \textbf{Replace with} & \textbf{Example pattern} \\
\midrule
ECDSA      & HIGH   & ML-DSA (FIPS 204)        & \texttt{ecrecover(}, \texttt{secp256k1} \\
RSA        & HIGH   & ML-KEM (FIPS 203)        & \texttt{RSA.}, \texttt{Crypto.PublicKey} import \\
DH         & HIGH   & ML-KEM (FIPS 203)        & \texttt{diffie\_hellman}, \texttt{ECDH} \\
SHA256     & MEDIUM & SHA-384/512 (FIPS 180-4) & \texttt{hashlib.sha256(}, \texttt{keccak256(} \\
AES\_SHORT & MEDIUM & AES-256-GCM (FIPS 197)   & \texttt{AES.new(...MODE\_}, \texttt{AES-128} \\
\bottomrule
\end{tabular}
\end{table}

A note on where patterns should start and stop, because both mistakes are easy
to make and we made both. Our first pattern list included a bare \texttt{dh}
followed by a dot, which matches any name ending in those letters, so ordinary
code containing \texttt{bandwidth} was reported as Diffie--Hellman. Requiring a
word boundary at the start fixes that. But requiring one at the \emph{end}
breaks matching in the other direction: \texttt{dilithium} then misses
\texttt{dilithium2}, and \texttt{diffie\_hellman} misses
\texttt{diffie\_hellman\_exchange}, since the next character is still a letter
or digit. Real library and function names nearly always have something
appended, so the rule must be lopsided: boundary at the start, never at the end.
The test set of Section~\ref{sec:benchmark} is what caught both faults.

\subsection{Hiding Comments and Text Without Losing Line Numbers}
\label{sec:masking}

Algorithm names turn up in comments and text strings far more often than people
expect, and it is worst in exactly the files we care about. A file that someone
is in the middle of migrating usually has comments explaining which old
algorithms were taken out. A tool that searches the raw text then reports those
comments as unfinished work, which is both wrong and very obvious to whoever
reads the report.

The obvious fix is to delete the comments and strings before searching. But that
shifts every character after them, so the line numbers come out wrong, and line
numbers are the most useful thing we produce.

So instead of deleting, we \emph{blank out}. Every character inside a comment or
a text string is replaced by a space, and line breaks are left alone. The result
is a copy of the file with exactly the same length and exactly the same lines,
but with all the prose turned into blank space. A position in this copy points
at the same place in the original file, so line numbers stay correct while the
search only ever sees real code.

The order of the patterns matters here. Triple-quoted strings and
\texttt{/* */} blocks have to be handled before single-line comments, otherwise
a \texttt{//} inside a documentation string cuts the search short. And the
patterns for single-line strings are not allowed to cross a line break, which
means an unmatched quote can spoil at most one line instead of swallowing the
rest of the file. This is pattern matching rather than proper parsing, and
Section~\ref{sec:validity} says where that falls short.

Names that appear only in the blanked-out parts are not simply dropped. We
report them in a separate list, which is not counted and does not affect the
score. This matters for goal G1: when a reviewer sees that the tool found
\texttt{ECDSA} in a comment and deliberately chose not to count it, they trust
the tool more. If we silently dropped it, they could not tell the difference
between a careful decision and a tool that just did not look.

\subsection{Counting One Place Once}
\label{sec:dedup}

Pattern lists like ours always overlap, because they have to catch both the
short and the long way of writing the same call. The line
\texttt{hashlib.sha256(payload)} matches both the full pattern and the plain
\texttt{sha256(} pattern. The import \texttt{from ecdsa import SigningKey,
SECP256k1} matches both the import pattern and the curve-name pattern. A naive
tool reports two results in each case, even though the developer only has one
place to go and fix.

So we group results that share both the same algorithm group and the same line
number, keeping the longest matched text because that describes the call most
precisely. We deliberately keep the algorithm group as part of the grouping key.
A line that really does use two different algorithms, such as an ECDSA signature
over a SHA-256 hash, should produce two results, because it needs two separate
fixes.

This is a correctness fix, not just tidying up, because the raw count feeds into
the risk score, the effort estimate and every total we print. On our test set
this step reduces 38 raw matches down to 26, and the gap between 26 and the 19
marked locations is entirely explained by the trick files, where the right
answer is zero.

\subsection{A Risk Score That Separates Files}
\label{sec:scoring}

Our earlier version scored a file by adding up severity points, 30 for each
serious finding, and capping the result at 100. The problem shows up
immediately: four serious findings already add up to 120, so any file with four
or more of them scores exactly 100 no matter how big it is or how many different
algorithms it uses. Across the scans we collected while building the tool, the
average score was 99.5, with almost no variation. A score that gives nearly
every file the same number tells you nothing, and since this number is printed
on compliance reports, that is a real problem rather than an academic one.

We replaced it with a combination of three things that vary independently of
each other:

\begin{equation}
S = \min\!\left(100,\;
\underbrace{\max_{f \in F} \sigma(f)}_{\text{how serious,} \le 60} +
\underbrace{\min(20,\, 5(|A| - 1))}_{\text{how many kinds,} \le 20} +
\underbrace{20\left(1 - e^{-\rho/5}\right)}_{\text{how dense,} \le 20}
\right)
\label{eq:score}
\end{equation}

\noindent Here $F$ is the set of findings and $\sigma$ turns a severity level
into points ($\text{HIGH} \mapsto 60$, $\text{MEDIUM} \mapsto 35$,
$\text{LOW} \mapsto 15$). $A$ is the set of different algorithm groups found,
and $\rho = 100|F| / L$ is the number of findings per hundred lines in a file of
$L$ lines.

Each part answers a different question: is anything completely broken, how many
different algorithms are involved, and how concentrated the problem is. The
third part uses an exponential curve so that a file packed with findings scores
higher than a sparse one, without one unusual file running away with the whole
scale. The score always goes up when any of the three inputs goes up, and only
a file that is both widespread and dense gets close to 100. We are not claiming
these weights are the best possible ones. They were chosen so the score is
understandable and tells files apart; Section~\ref{sec:validity} notes that
checking them against expert opinion is still to be done.

% ============================================================================
\section{Implementation}
\label{sec:impl}

\subsection{The Main Version}

The main version is written in Python and exposed through a FastAPI web service
with twelve endpoints, covering scanning, migration planning, report
generation, the CycloneDX export, certificate issuing and blockchain
recording. Results are stored as files on disk rather than in a database, which
keeps the setup simple at the cost of handling many users at once. For a tool
that is run once per file, we think that is the right trade.

Two bugs we found while hardening it are worth writing down, because both are
common in research prototypes. First, the blockchain client was created when
the program started, so a missing signing key stopped the whole service from
running instead of just disabling one feature; it is now created only when
first needed, and only the blockchain endpoints fail without it. Second, the
endpoint that writes to the blockchain, which spends real funds from a key held
on the server, had no authentication at all, so anyone who found the URL could
drain the operator's balance. The design in Section~\ref{sec:anchor} removes
that danger completely rather than just putting a lock on it.

\subsection{The Same Tool Twice, Checked to Match}
\label{sec:dualruntime}

Goal G3 says a company should be able to check its private code without sending
it anywhere. Since our detection is just pattern matching on text, it does not
need a server at all, so we wrote the whole detection path a second time in
JavaScript. This version reads the file using the browser's own file API, does
the blanking, searching, grouping and scoring inside the page, and builds the
reports and the CycloneDX list there too. Nothing is uploaded.

Keeping two copies of the same logic invites them to drift apart, which would be
a serious problem: the version we measured and the version people actually use
would no longer be the same tool. So we treat matching as something to check
rather than hope for. A test script runs the JavaScript detection code outside
the browser and compares it against the Python version on all seven test files,
checking six things: the set of findings as (line, algorithm) pairs, the number
of findings, the list of comment-only mentions, which post-quantum algorithms
were spotted, the risk score and the line count. All seven files agree on all
six. The check is automatic and is meant to be run whenever the patterns change.

\subsection{The Machine-Readable Inventory}
\label{sec:cbom}

Every scan can be exported as a CycloneDX 1.6~\cite{cyclonedx} document, where
each algorithm group found becomes one \texttt{cryptographic-asset} entry
carrying its kind, its severity, how many times it appears, which lines it is
on, the recommended replacement and the standard that governs it. Using the
standard format instead of our own means the output can go straight into the
supply-chain review tools a company already runs, and lets them track their
cryptography from one release to the next.

\subsection{Recording the Result Without Holding Any Keys}
\label{sec:anchor}

Goal G4 says someone else must be able to check a report is genuine. We take the
scan's details --- its identifier, the company, the time, the number of findings
and the score --- compute a Keccak-256 hash of them, and write that hash to a
small registry contract on the Ethereum Sepolia test network. The blockchain
supplies the two things we need: proof the record has not been altered, and
proof of when it was made. Anyone holding the report can recompute the hash and
check it matches, without having to trust whoever ran the scan. Only the hash
goes on-chain, so no source code is revealed.

How the transaction gets signed matters more than it first appears. With a key
held on the server, as in our earlier version, the operator becomes a custodian:
the key has to be funded, protected and rotated, and every endpoint that uses it
is a way to drain the balance. Instead, we have the transaction signed by the
wallet of whoever is running the scan, in their own browser. The service holds
no key, so there is nothing to steal and nothing to drain, and the record is
tied to the party that actually did the scan. That is stronger proof of who did
what than a single shared service key could give.

\section{Evaluation}
\label{sec:eval}

\subsection{How We Built the Test Set}
\label{sec:benchmark}

We built a test set of seven source files and marked by hand every place where
quantum-weak cryptography is actually used. Each mark records a line number and
an algorithm group. The tool only gets credit for a result if both of these
match. We chose this strict rule on purpose: telling a developer the right
algorithm but the wrong line does not help them do the work.

We wrote the marking rules down before marking anything, and we release them
with the data. Imports count, because importing a weak algorithm is part of what
has to be migrated. Function calls count. A function definition counts if its
body implements the algorithm. Mentions inside comments, documentation strings
and text strings do not count. Names that merely contain an algorithm name
inside them do not count either.

Four of the files contain real problems: a Solidity consent registry using
\texttt{ecrecover} and \texttt{keccak256}; a Python settlement signer using
ECDSA, RSA, SHA-256, short-key AES and a Diffie--Hellman exchange; a Python key
service; and a half-migrated Solidity contract where an ML-DSA path has been
added but the old \texttt{ecrecover} path is still live. That last one matters
because half-migrated code is the normal real-world state, and the tool has to
still report what is left.

The other three files contain nothing to find, and they are there to try to
trick the tool. The first is fully post-quantum code importing
\texttt{dilithium2} and \texttt{kyber512} and hashing with SHA-384. The second
describes a finished migration in prose, mentioning \texttt{RSA.generate},
\texttt{hashlib.sha256}, \texttt{ECDSA.sign}, \texttt{Diffie-Hellman} and
\texttt{ecrecover} across comments, a documentation string and a text string,
while calling none of them. The third contains names like \texttt{bandwidth},
\texttt{sha256\_cache\_enabled}, \texttt{rsa\_config\_path} and
\texttt{ecdsa\_feature\_flag}, none of which has anything to do with
cryptography. Each of the three targets one of the three mistakes described in
Section~\ref{sec:method}. In total the test set has nineteen marked locations.

\subsection{The Three Versions We Compare}

To show which fix does what, we run the tool in three versions:

\begin{itemize}
  \item \textbf{v1-baseline}: our original pattern list, searching the raw text,
  with no grouping. This is a real earlier version of the tool, kept exactly as
  it was, not something we rebuilt afterwards to look bad.
  \item \textbf{v2-boundaries}: the corrected pattern list with word boundaries,
  still searching raw text with no grouping.
  \item \textbf{v3-full}: the corrected patterns plus comment blanking plus
  per-line grouping. This is the version we ship.
\end{itemize}

\subsection{Results}
\label{sec:results}

\begin{table}[t]
\centering
\caption{The three versions compared. Seven files, nineteen marked locations.
``Raw hits'' is the total number of pattern matches before grouping.}
\label{tab:ablation}
\begin{tabular}{@{}lrrrrrrr@{}}
\toprule
\textbf{Version} & \textbf{TP} & \textbf{FP} & \textbf{FN} &
\textbf{Prec.} & \textbf{Rec.} & \textbf{F1} & \textbf{Raw hits} \\
\midrule
v1-baseline   & 19 & 11 & 0 & 0.633 & 1.000 & 0.776 & 39 \\
v2-boundaries & 19 & 11 & 0 & 0.633 & 1.000 & 0.776 & 38 \\
v3-full       & 19 &  0 & 0 & \textbf{1.000} & 1.000 & \textbf{1.000} & 26 \\
\bottomrule
\end{tabular}
\end{table}

Table~\ref{tab:ablation} shows the comparison. Three things stand out.

First, recall stays at 1.000 the whole way through. Our fixes never hide a real
problem. This is the thing that had to be true for them to be worth having: in a
security tool, buying accuracy by missing real cases would be a bad trade.

Second, adding word boundaries on its own did not help at all. Version v2
removes one raw match compared with v1, but it does not remove a single false
positive. This is worth reporting rather than hiding. The mistake that word
boundaries fix is real, and we wrote a file specifically to trigger it, but on
this data almost all the damage comes from matching comments, which word
boundaries do nothing about. We mention it because adding word boundaries is the
obvious first thing to try, and on its own it is not enough.

Third, blanking comments together with per-line grouping removes all eleven
false positives. Those eleven are matches inside comments and documentation
strings, in both the real files and the trick files, plus duplicate matches on
the same line. Broken down by algorithm, precision and recall are both 1.000 for
ECDSA (4 locations), RSA (4), SHA-256 (8), Diffie--Hellman (2) and short-key AES
(1).

\paragraph{What these numbers do and do not prove.} The scores should not be
read as how accurate the tool is on real code, and the reason is simple. The
test set is small, at seven files and nineteen locations, and we wrote it
ourselves while building the tool. Under those conditions, a perfect score for
our final version is close to being guaranteed by construction: we wrote the
trick files to catch mistakes we had already decided to fix. What the table does
show is narrower but still worth having: each fix removes the mistake it was
meant to remove, it removes it completely even on data designed to be awkward,
and nothing real is lost in the process. The number to look at is the
\emph{change}, from 0.633 to 1.000, under a comparison that shows which fix
caused it. Getting a real accuracy figure needs a test set we did not write
ourselves, which is the main piece of future work in Section~\ref{sec:future}.

The small size makes this worse in a second way. Even with all nineteen correct,
nineteen results only support a 95\% confidence range for precision of
$[0.832, 1.000]$. Per algorithm it is far weaker: the single short-key AES
location gives $[0.207, 1.000]$, which tells us essentially nothing about that
algorithm. So we do not claim a particular level of accuracy from this test set,
only a difference between the three versions measured on the same data. To
narrow the range enough to claim accuracy above 0.95, we would need roughly 150
marked locations, which is what sets the size of the test set described in
Section~\ref{sec:future}.

\subsection{Does the New Score Tell Files Apart?}

\begin{table}[t]
\centering
\caption{Risk scores for every test file under both scoring methods. The old
method gives the two worst files exactly the same number; the new one separates
all four.}
\label{tab:scores}
\begin{tabular}{@{}lrrrrr@{}}
\toprule
\textbf{File} & \textbf{Found} & \textbf{Kinds} & \textbf{Lines} &
\textbf{Old score} & \textbf{New score} \\
\midrule
Python settlement signer  & 9 & 4 & 108 & 100.0 & 91.2 \\
Python key service        & 5 & 4 &  26 & 100.0 & 94.6 \\
Solidity consent registry & 3 & 2 &  81 &  60.0 & 75.5 \\
Half-migrated contract    & 2 & 2 &  32 &  45.0 & 79.3 \\
Post-quantum module       & 0 & 0 &  39 &   0.0 &  0.0 \\
Prose-only module         & 0 & 0 &  23 &   0.0 &  0.0 \\
Trick-names module        & 0 & 0 &  38 &   0.0 &  0.0 \\
\bottomrule
\end{tabular}
\end{table}

Table~\ref{tab:scores} compares the two scoring methods on every test file. The
problem we predicted in Section~\ref{sec:scoring} shows up exactly where
expected: the two files whose severity points go past the cap both get 100.0, so
the old method cannot tell apart a file with 9 findings in 108 lines from one
with 5 findings in 26 lines, even though the second is more than twice as
densely affected. The 19 scans we collected while building the tool show the
same thing on a larger scale, with an average of 99.5 and almost no variation,
because real files routinely have more than four serious findings. The new
method gives four different numbers to the four affected files, and correctly
gives zero to all three clean ones.

Two things about the new score are worth pointing out rather than hiding. First,
on small files it can give a \emph{higher} number than the old one: the
half-migrated contract goes up from 45.0 to 79.3, because two live serious
problems inside 32 lines really is dense exposure, which simply counting
findings understates. Second, the order is not the order of finding counts. The
key service scores above the settlement signer despite having fewer findings,
because the density part dominates: 19.2 findings per hundred lines against 8.3.
This is what we intended --- a small file that is almost entirely affected needs
attention sooner than a big file with a few scattered calls --- but it does mean
the score should be read as how concentrated the problem is, not as a stand-in
for how many findings there are.

We should be clear about what this does and does not show. It shows the new
score tells files apart where the old one could not, and that its behaviour
follows from the three parts we defined rather than from hitting a cap. It does
not show that its ordering matches what a security expert would say, which would
need a ranking study we have not done. Section~\ref{sec:future} lists this.

\subsection{Do Both Versions Agree?}

The checking script reports that the Python and JavaScript versions agree on all
seven test files, across all six things it compares. So the browser version that
people actually use behaves exactly like the version we measured, rather than
approximately like it. We think this is a necessary condition for reporting
results from one version while shipping another.

% ============================================================================
\section{Discussion}
\label{sec:discussion}

\subsection{Fitting Existing Regulations}

What the tool produces lines up with obligations Indian companies already have.
The RBI cybersecurity framework requires regulated companies to know what
cryptography they are using and to document how they are fixing problems.
CERT-In's directions require reporting on assets and incidents, which you cannot
do without knowing what cryptography you have. Internationally, NIST's
timetable~\cite{nistir8547} sets dates by which the old algorithms are no longer
allowed, so companies need to be able to show progress against it, and ISO/IEC
27001 requires a written cryptography policy and evidence that it is followed.

Together, the dated blockchain record and the standard-format inventory solve
one specific awkward problem: proving \emph{when} a check was done to somebody
who was not there at the time. The blockchain record does this without the other
party having to trust the company, and without revealing any source code.

\subsection{Which Industries Should Worry First}

How much risk a company carries depends on what it uses its blockchain for.
Pharmaceutical and manufacturing companies track batches through a supply chain,
where signatures prove custody, so a forged old signature corrupts a record
regulators depend on. Healthcare anchors consent and patient records, where the
backwards-looking threat is worst because the data must stay confidential for a
lifetime. Financial companies have shorter confidentiality needs but more money
at immediate risk. By Mosca's rule~\cite{mosca2018} the first two are more
urgent than the attention they currently get, because their secrets must last
longest.

\subsection{What the Tool Cannot Do}
\label{sec:validity}

We separate the limits of the method from the limits of our testing.

\paragraph{The method.} Our detection works on text. It finds where a
cryptography library is called, but it cannot tell whether a correctly-named
algorithm is being used correctly, and it cannot follow calls between functions.
If a project wraps ECDSA inside its own \texttt{sign()} function and then calls
that from twenty places, we report one result at the wrapper, not twenty. For
planning a migration that is often what you want, since the wrapper is the right
place to change, but it does understate how much of the system is affected.
Algorithms chosen at run time, cryptography coming from a dependency, and
compiled code are all outside what we handle. The comment blanking is pattern
matching rather than real parsing, so unusual or nested string delimiters can
confuse it; we limit the damage by not letting single-line string patterns cross
a line break, but a proper parser for each language would be strictly better.

\paragraph{Our testing.} Three problems apply, in order of how serious they are.
\emph{We wrote the test set}: the trick files target mistakes we had already
decided to fix, so the v3 scores are optimistic. This is the biggest weakness of
the evaluation and we are not going to play it down. \emph{The test set is small
and artificial}: seven files of a few dozen lines each cannot stand in for real
blockchain projects, which are bigger, use more libraries, and contain patterns
we have not thought of. \emph{The matching rule is harsh}: requiring both the
line and the algorithm to match exactly is defensible, but a result one line
away from the mark counts as both a false positive and a miss, which punishes
near-misses harder than a developer would.

\paragraph{The rest of the system.} The blockchain records go on a test network,
so they have no real economic weight; a production deployment would use the main
network or a private one. Certificates are tracked in a local file rather than a
public service, so right now a certificate is evidence that a check was done
rather than something anyone can independently look up.

\subsection{Could This Tool Help an Attacker?}

A tool that points out weak cryptography can obviously be misused. But the
weakness it reports is not secret: that ECDSA falls to Shor's algorithm is
public knowledge, and the keys in question are already on public blockchains, so
an attacker learns nothing they could not get directly. The balance favours
defenders, who must find every place, over attackers, who need only one. Two
design choices follow: the browser version means a company never hands its
source to anyone, and only hashes go on-chain, never code or findings. The tool
also refuses to issue a certificate while any finding remains, and its reports
state the method's limits inside the document.

\section{Future Work}
\label{sec:future}

The most important thing to do next is to test the tool on code we did not
write, and we have already built and released the machinery for it so that the
method is fixed in advance rather than decided afterwards once we have seen the
results. It works in three steps, and keeping the steps apart is what makes it
trustworthy.

Step one picks random recent blocks, finds the transactions that created
contracts, keeps those with published source, and produces a list of
\emph{candidate lines} for a person to judge. The search that picks candidates is
deliberately worse than our detector: plain text fragments, no word boundaries,
no comment blanking, no grouping. That makes the candidate list a strict
superset of everything our detector could report, which is what lets us measure
the cases it \emph{misses} as well as the ones it gets wrong. Step three checks
that this superset property actually holds.

Step two shows a person each candidate line with its surrounding lines and asks
what that line really does. The tool does not load our detector and never shows
its prediction. This is the crucial part: somebody who can see a prediction
tends to agree with it, so the labels would copy the detector's mistakes and the
accuracy figure would be meaningless. Two people label independently, and where
they disagree we drop the line rather than taking a majority vote, which would
invent confidence the data does not support.

Step three reports agreement between the two people using Cohen's $\kappa$,
along with precision, recall and $F_1$ per algorithm with confidence intervals.
We treat $\kappa$ as a pass-or-fail gate. Below about 0.6 means our marking
rules are unclear, and the right response is to rewrite them and mark again, not
to publish the accuracy figures anyway.

Three smaller things follow. A proper parser for each language would remove the
text-matching approximation of Section~\ref{sec:validity} and let us follow
calls between functions, so wrappers no longer hide the true size of the
problem. Checking the weights in Eq.~\ref{eq:score} against how experts rank
files would turn the score from something that separates files into something
validated. And measuring the real gas cost of on-chain lattice verification
would tell us when the off-chain workaround can be dropped.

% ============================================================================
\section{Conclusion}
\label{sec:conclusion}

Companies running blockchain systems must migrate their cryptography, and the
first thing they need is a list of where the weak algorithms are used. Their
existing tools do not produce it: smart contract checkers look for logic bugs,
and cryptography checkers look for algorithms used wrongly, whereas the problem
here is algorithms used perfectly correctly that are about to stop being safe.
We built a tool that produces that list and carries it through to a migration
plan, a standard-format inventory, a regulator's report, and a verifiable
record.

Our real finding is about accuracy. A tool that searches code for cryptography
is easy to make find everything and hard to stop from being wrong, and its
mistakes are not random: two causes account for nearly all of them. Names
matched inside comments and text strings, and one line counted several times
because several patterns matched it. Blanking comments without losing line
numbers, and grouping matches per line, remove both, taking precision from
0.633 to 1.000 without losing a single real case. Adding word boundaries alone,
the obvious first thing to try, removes no false positives at all. We also
showed that the severity-adding score such tools normally use saturates, giving
the same maximum to files differing several times over in how affected they are,
and gave a three-part replacement that tells them apart.

We are deliberate about what we have not shown. Because we wrote the test set
ourselves while building the tool, the final numbers are optimistic, and the
result we stand behind is the change from 0.633 to 1.000 rather than the
endpoint. Nineteen cases cannot support an accuracy claim in any case: the
confidence interval is too wide. Establishing real accuracy needs a test set
labelled by people who did not build the tool, which is where we go next. The
tool, the test set, the marking rules and the measurement scripts are all
released, so both the result and its limits can be checked.

% ----------------------------------------------------------------------------
% Fill in or delete this whole section before submitting.
\subsubsection*{Acknowledgements.}
% <ACKNOWLEDGEMENTS: funding, supervisor, institution -- or delete the section>

\subsubsection*{Disclosure of Interests.}
The author has no competing interests to declare.

% ============================================================================
\begin{thebibliography}{99}

\bibitem{shor1997}
Shor, P.W.: Polynomial-time algorithms for prime factorization and discrete
logarithms on a quantum computer. SIAM J. Comput. \textbf{26}(5), 1484--1509
(1997)

\bibitem{grover1996}
Grover, L.K.: A fast quantum mechanical algorithm for database search. In: Proc.
28th Annual ACM Symposium on Theory of Computing (STOC), pp. 212--219. ACM
(1996)

\bibitem{nakamoto2008}
Nakamoto, S.: Bitcoin: a peer-to-peer electronic cash system (2008)

\bibitem{fips203}
National Institute of Standards and Technology: Module-Lattice-Based
Key-Encapsulation Mechanism Standard. FIPS 203, NIST, Gaithersburg, MD (2024)

\bibitem{fips204}
National Institute of Standards and Technology: Module-Lattice-Based Digital
Signature Standard. FIPS 204, NIST, Gaithersburg, MD (2024)

\bibitem{fips205}
National Institute of Standards and Technology: Stateless Hash-Based Digital
Signature Standard. FIPS 205, NIST, Gaithersburg, MD (2024)

\bibitem{nistir8547}
National Institute of Standards and Technology: Transition to Post-Quantum
Cryptography Standards. NIST IR 8547, NIST, Gaithersburg, MD (2024)

\bibitem{bernstein2017}
Bernstein, D.J., Lange, T.: Post-quantum cryptography. Nature
\textbf{549}(7671), 188--194 (2017)

\bibitem{mosca2018}
Mosca, M.: Cybersecurity in an era with quantum computers: will we be ready?
IEEE Secur. Priv. \textbf{16}(5), 38--41 (2018)

\bibitem{webber2022}
Webber, M., Elfving, V., Weidt, S., Hensinger, W.B.: The impact of hardware
specifications on reaching quantum advantage in the fault tolerant regime. AVS
Quantum Sci. \textbf{4}(1), 013801 (2022)

\bibitem{slither}
Feist, J., Grieco, G., Groce, A.: Slither: a static analysis framework for smart
contracts. In: Proc. 2nd Int. Workshop on Emerging Trends in Software
Engineering for Blockchain (WETSEB), pp. 8--15. IEEE (2019)

\bibitem{cryptoguard}
Rahaman, S., Xiao, Y., Afrose, S., Shaon, F., Tian, K., Frantz, M.,
Kantarcioglu, M., Yao, D.: CryptoGuard: high precision detection of
cryptographic vulnerabilities in massive-sized Java projects. In: Proc. ACM
SIGSAC Conf. on Computer and Communications Security (CCS), pp. 2455--2472. ACM
(2019)

\bibitem{crysl}
Kr\"uger, S., Sp\"ath, J., Ali, K., Bodden, E., Mezini, M.: CrySL: an extensible
approach to validating the correct usage of cryptographic APIs. IEEE Trans.
Softw. Eng. \textbf{47}(11), 2382--2400 (2021)

\bibitem{cyclonedx}
OWASP Foundation: CycloneDX Bill of Materials Specification, version 1.6 (2024)

\bibitem{sp80057}
Barker, E.: Recommendation for Key Management: Part 1 --- General. NIST SP
800-57 Part 1 Rev. 5, NIST, Gaithersburg, MD (2020)

\end{thebibliography}

\end{document}
